\section{Collision vs Trigger}

Collision dan trigger menjadi landasan mekanik game. Bab ini membahas OnTriggerEnter dan OnCollisionEnter. \class{Collider} mendeteksi benturan \cite{unityPhysics}. Dua mode:

\begin{itemize}
  \item \textbf{Collision}: Objek saling bertabrakan secara fisik (memantul, terdorong). Untuk dinding, lantai, objek padat
  \item \textbf{Trigger}: Hanya mendeteksi overlap, tidak ada respon fisik. Centang ``Is Trigger'' di Inspector. Untuk collectibles, checkpoint, zona
\end{itemize}

Untuk collectibles: gunakan \textbf{Trigger} agar bola bisa ``melewati'' dan mengumpulkan tanpa mendorong objek. Jika pakai Collision biasa, bola akan mendorong collectible.


\subsection{OnTriggerEnter vs OnCollisionEnter}

\begin{itemize}
  \item \texttt{OnTriggerEnter(Collider other)}: Dipanggil saat collider lain masuk trigger. Gunakan untuk collectibles
  \item \texttt{OnCollisionEnter(Collision collision)}: Dipanggil saat collision fisik terjadi. Gunakan untuk tabrakan dengan dinding
\end{itemize}

\subsection{OnTriggerEnter}

\begin{unitycode}
void OnTriggerEnter(Collider other)
{
    if (other.gameObject.CompareTag("PickUp"))
    {
        other.gameObject.SetActive(false);
    }
}
\end{unitycode}
